\par\Needspace{4\baselineskip}
\begin{center}
\SourceLeaf{040}\pdfbookmark[2]{22.}{article-022}22.
\end{center}
Durch Differenzieren von Formel IV erhält man (wenn man $\psi$ als konstant betrachtet)
\[
\begin{aligned}
\frac{\mathrm{d}u}{u}&=\frac{1}{2}\left(\operatorname{tang}\frac{1}{2}(v+\psi)-\operatorname{tang}\frac{1}{2}(v-\psi)\right)\mathrm{d}v\\
&=\frac{r\operatorname{tang}\psi}{p}\,\mathrm{d}v;
\end{aligned}
\]
daraus folgt
\[
r r\,\mathrm{d}v=\frac{p r}{u\operatorname{tang}\psi}\,\mathrm{d}u
\]
oder, wenn man für $r$ den Wert aus X einsetzt,
\[
r r\,\mathrm{d}v=b b\operatorname{tang}\psi\left\{\frac{1}{2}e\left(1+\frac{1}{u u}\right)-\frac{1}{u}\right\}\mathrm{d}u.
\]
Integriert man dann so, dass das Integral im Perihel verschwindet, so ergibt sich
\[
\begin{aligned}
\int r r\,\mathrm{d}v&=b b\operatorname{tang}\psi\left\{\frac{1}{2}e\left(u-\frac{1}{u}\right)-\log u\right\}\\
&=k t\BareRoot{p}\, .\,\BareRoot{(1+\mu)}=k t\operatorname{tang}\psi\, .\,\BareRoot{b}\, .\,\BareRoot{(1+\mu)}.
\end{aligned}
\]
Der Logarithmus ist hier hyperbolisch. Verwendet man Logarithmen des \textsc{Briggs}schen Systems oder allgemeiner eines Systems mit dem Modul $=\lambda$ und vernachlässigt man die Masse $\mu$ (von der wir annehmen dürfen, dass sie sich bei einem Körper auf einer hyperbolischen Bahn nicht bestimmen lässt), so erhält die Gleichung die folgende Form.
\[
\mathrm{XI.}\quad\begin{gathered}
\frac{1}{2}\lambda e\frac{u u-1}{u}-\log u=\frac{\lambda k t}{b^{\frac{3}{2}}},\\
\text{oder mit }F\\
\lambda e\operatorname{tang}F-\log\operatorname{tang}(45^0+\frac{1}{2}F)=\frac{\lambda k t}{b^{\frac{3}{2}}}\cdot
\end{gathered}
\]
Gehen wir von \textsc{Briggs}schen Logarithmen aus, so haben wir \\
$\log\lambda=9{,}637\,7843\,113$, $\log\lambda k=7{,}873\,3657\,527$. Eine etwas höhere Genauigkeit lässt sich jedoch erzielen, wenn man unmittelbar hyperbolische Logarithmen verwendet. Die hyperbolischen Logarithmen der Tangenswerte finden sich in mehreren Tafelsammlungen, etwa in der von \textsc{Schulze} herausgegebenen Sammlung, und für einen noch größeren Bereich im Magnus Canon Triangulorum Logarithmicus von \textsc{Beni. Ursinus}, Colon. 1624, wo sie in Schritten von $10^{\prime\prime}$ angegeben sind. — Aus Formel XI geht ferner hervor, dass reziproken Werten von $u$ oder entgegengesetzten Werten von $F$ und $v$ entgegengesetzte Werte von $t$ entsprechen. Daher werden gleich große Abschnitte der Hyperbel, die auf beiden Seiten des Perihels gleich weit von ihm entfernt liegen, in gleichen Zeiten durchlaufen.
